% !TEX root = ../main.tex
\section{Introduction}
\label{sec:introduction}

\subsection{Problem Background}

埃博拉病（Ebola virus disease, EVD）由埃博拉病毒引起，主要经直接接触感染者或病亡者的体液传播，病死率高。1976 年该病在扎伊尔（今刚果民主共和国）被首次识别，此后非洲已报告二十余次暴发；其中 2014--2016 年西非疫情的规模最大，报告病例超过 28\,000 例、死亡超过 11\,000 人\cite{whofactsheet,reliefweb2016}。

2026 年 5 月，刚果民主共和国再次暴发疫情。到 2026 年 10 月 25--26 日，本轮疫情已持续约五个月，累计确诊约 10\,060 例、累计死亡约 4\,887 人，粗病死率约 48.6\%，波及 8 个省的 66 个卫生区\cite{whodon622,ebolatrader}。本轮由本迪布焦病毒引起，与常见的扎伊尔毒株不同，目前尚无获批的疫苗或特异性治疗药物；世界卫生组织已在 5 月宣布本轮疫情构成国际关注的突发公共卫生事件\cite{whosituation2026}。

从公开通报看，防控的困难并不在单一环节。八成以上的新发病例是在已登记的接触者之外被发现的，说明现有追踪没有覆盖真实的传播面；约三分之二的死亡发生在社区，患者生前未在医疗机构就诊，说明葬礼和家庭内接触仍在传播中占有主要份额，不能只当作院内传播的补充。伊图里、北基伍等地武装冲突持续，治疗中心曾遭袭击，检测实验室从 1 家增加到 16 家、收治床位超过 800 张、接触者随访率接近 80\% 之后，世界卫生组织与非洲疾病预防控制中心的联合应对计划仍有超过 4 亿美元的资金缺口\cite{whodg2026}。境外方面，肯尼亚已报告 9 例确诊病例，尼日利亚于 10 月 20 日确诊首例，该病例有乌干达旅行史\cite{whodon622}。

\subsection{Problem Restatement and Task Breakdown}

题目要求围绕本轮疫情完成六件事。

\begin{enumerate}[label=\arabic*.,leftmargin=1.8em]
  \item 建立描述传播动力学的连续时间模型，逐条说明假设及其生物学或流行病学依据，给出完整常微分方程组并解释每一状态转移的机制与参数含义；模型要包含人际、院内和葬礼三条传播途径，并用时变参数反映干预强度的变化。
  \item 收集疫情数据，说明来源与预处理方法，估计关键参数，评价拟合优度，并讨论不确定性从何而来。
  \item 在第一个任务建立的模型上，逐项考察患者隔离、接触者追踪、安全安葬、院内感染控制、疫苗接种、MBP134/瑞德西韦治疗、奥贝德西韦暴露后预防和社区健康教育对累计病例、死亡人数与疫情持续时间的影响；比较若干组合方案；比较 7 月中旬（2\,073 例）即加强干预与现阶段（10\,060 例）才加强干预的差别；测算把随访率从 78\% 提高到 90\% 或 95\% 需要多少资源、又能减少多少传播。
  \item 建立空间传播模型，考虑省际人口流动矩阵、出行限制与边境管控、区域联防联控和环围疫苗接种，分析各种空间策略对扩散的影响，判断局部暴发在什么条件下会演变为大规模流行，并评估肯尼亚的 9 例输入和尼日利亚的首例是否可能引起持续的本地传播、应当采取什么措施。
  \item 做敏感性分析，找出对结果影响最大的参数，讨论数据质量、参数估计、毒株特性、冲突局势与人群行为变化带来的不确定性，检验模型在不同情景下的稳健性并说明局限。
  \item 面向刚果民主共和国卫生部与国际捐助方撰写执行摘要，交代背景，给出核心发现、最紧迫的行动建议、资源需求与资金缺口，并说明为什么这属于国家安全投资而不只是慈善。
\end{enumerate}

提交的 PDF 不超过 25 页，依次为单页摘要表、目录、正文（执行摘要可以并入）与参考文献；若使用人工智能，另附使用报告，不计入页数。

\subsection{Our Work}

本文的分析沿一条线推进：先把多层级通报数据整理成可以用于估计的周序列，在八舱连续时间模型上估计时变传播强度与检出比例，再把标定好的模型分别用于干预情景、空间与跨境、敏感性三个方向，最后汇总成执行摘要（图~\ref{fig:workflow}）。

\begin{figure}[htbp]
\centering
\begin{tikzpicture}[
  node distance=5mm and 6mm,
  box/.style={draw,rounded corners=2pt,align=center,font=\small,inner sep=4pt,minimum height=9mm},
  hub/.style={box,fill=black!6},
  arr/.style={-{Latex[length=1.6mm]},thick}
]
\node[box] (data) {通报数据\\三层口径分离};
\node[hub,right=of data] (model) {八舱连续\\时间模型};
\node[box,right=of model] (est) {样条 $k(t)$\\与 $\rho$ 估计};
\node[box,right=of est] (int) {干预情景\\单措施/组合/时机};
\node[box,below=8mm of int] (spa) {空间与跨境\\按 $R_{\mathrm{local}}$ 分情景};
\node[box,left=of spa] (sen) {敏感性\\参数包络};
\node[box,left=of sen] (pol) {执行摘要\\优先级与资金};
\draw[arr] (data) -- (model);
\draw[arr] (model) -- (est);
\draw[arr] (est) -- (int);
\draw[arr] (int) -- (spa);
\draw[arr] (spa) -- (sen);
\draw[arr] (sen) -- (pol);
\end{tikzpicture}
\caption{建模与交付的先后关系}
\label{fig:workflow}
\end{figure}

模型把社区、院内、葬礼三条通道写进同一个感染力，省与家庭两层索引分别承担人口流动与葬礼混合池的作用。这样处理的理由很直接：三条通道若各用一套参数，院内与葬礼传播的相对强弱就无从比较，"加快收治反而增加病例"这类结果在结构上根本不会出现。

参数估计使用 27 周全国序列，第 10 周的补报单独处理。时变强度取 16 节点对数样条加二阶差分惩罚，与检出比例一同估计，得到样本内 $R^2=0.955$、RMSE 为 34.6 例/周；用前 20 周拟合、后 7 周回测，RMSE 为 119.5、MAPE 为 20.2\%。检出比例估计为 0.392，后期有效再生数为 0.893。全文只保留这一套参数，摘要、正文和政策部分的数字都由它生成。

由此得到的主要结果如下。从第 27 周外推 26 周，基线新增 5\,672 例、报告口径死亡 2\,937 例；单措施中安全安葬最有效，累计病例降至基线的 $61.7\%$ 以下（30 周内收尾），接种 5\%/周与治疗次之，院内感染控制为 $-17.1\%$；四件套组合 C2 可减少 $67.7\%$，加入接种与治疗的 C3 减少 $72.0\%$。干预提前到 7 月中旬时，同期末病例由 17\,188 例降到 6\,884 例。有两处结果与直觉相反：单独加快隔离、单独提高追踪率会让 26 周新增分别增加 2.0\% 与 3.4\%，原因是标定得到的院内人均感染力约为社区的两倍，收治加快等于把病人更快送进更危险的舱室；出行限制改变的是空间分布和响应时间，压不下全国总量。前一点在组合情景中得到了反证，配上院内感染控制之后，同样的动作变为 $-20.0\%$。

不确定性没有压成单点数字。有效再生数同时给出条件统计区间 $[0.86,0.93]$ 与强度整体扰动 $\pm30\%$ 后的包络 $[0.62,1.22]$，后者跨过 1，因此下文不把"已经低于 1"当作既定前提。跨境风险也没有把刚果（金）的估计值直接推给邻国，而是按当地再生数取 0.8、1.0、1.2、1.5 四种情景分别判断：不超过 1 时输入链必然熄灭，1.2 时局部暴发概率为 19\%--23\%，1.5 时降至 3\%--4\%。
